\section*{Chapter \ref{ch:in:quantitative-hedge-funds-and-systematic-managers} --- Quantitative Hedge Funds and Systematic Managers}

\begin{solution}{exo:in:quantitative-hedge-funds-and-systematic-managers:1}
$92.0/300=\$307$ million of RAUM per employee; $160/300=53\%$ of staff in advisory functions.
\end{solution}

\begin{solution}{exo:in:quantitative-hedge-funds-and-systematic-managers:2}
About \$7 billion: RAUM counts the securities portfolio at gross value and does not deduct the fund's borrowing or
other liabilities, so gross positions, not net assets, are reported.
\end{solution}

\begin{solution}{exo:in:quantitative-hedge-funds-and-systematic-managers:3}
No: an adviser solely to qualifying private funds with less than \$150 million of private-fund assets in the United
States is exempt from registration. It is an \omterm{def:in:quantitative-hedge-funds-and-systematic-managers:era}{exempt reporting adviser} and files a report on Form ADV.
\end{solution}

\begin{solution}{exo:in:quantitative-hedge-funds-and-systematic-managers:4}
Only Winton Capital Management ($4.6/10.1=0.46$); AQR is close ($87.6/160.5=0.55$). Much of their RAUM sits outside
private funds: managed accounts or registered funds, which a hedge-fund count alone would miss.
\end{solution}

\begin{solution}{exo:in:quantitative-hedge-funds-and-systematic-managers:5}
Aspect Capital, \$24.4 million per employee, is at about the 9th percentile; Squarepoint's US adviser, \$2.57 billion, at
about the 99.6th. The latter is a registered entity of 71 people through which assets are advised while the firm's
staff work for affiliates elsewhere: the ratio divides a group's assets by an entity's headcount.
\end{solution}

\begin{solution}{exo:in:quantitative-hedge-funds-and-systematic-managers:6}
The bins at \$1 billion and above hold $18+4+3+0+1=26$ of 1\,457 advisers: 1.8\%.
\end{solution}

\begin{solution}{exo:in:quantitative-hedge-funds-and-systematic-managers:7}
Ten equal advisers: top-ten share 100\%, index 1\,000. A thousand equal: 1\% and 10. An index of 78 is what about
$10\,000/78=128$ equal-sized advisers would give: a large, unconcentrated industry with some big firms.
\end{solution}

\begin{solution}{exo:in:quantitative-hedge-funds-and-systematic-managers:8}
The row is one registered adviser; the group's thousands of staff are employed by affiliates that are not in the row.
RAUM per employee of an entity is not the firm's productivity, and RAUM is gross assets, not revenue or profit.
\end{solution}

\begin{solution}{pb:in:quantitative-hedge-funds-and-systematic-managers:1}
\begin{enumerate}
\item Systematic: trades decided by rules in software; discretionary: by portfolio managers. Research-driven
multi-strategy funds; trend followers; statistical-arbitrage houses.
\item 5.A employees; 5.B(1) advisory employees; 5.F RAUM; 7.B private funds.
\item Securities portfolios under continuous supervision at gross value, with no deduction of indebtedness; AUM as
reported to investors is usually net.
\item An adviser exempt from registration that files a short Form ADV report; the private-fund exemption applies below
\$150 million of US private-fund assets.
\item One adviser is not the group; no returns, fees, pay or revenue; answers may be most of a year old.
\item 2\,577; 1\,458.
\item \$119 million; \$58 million and \$221 million.
\item 37\% and 62\% (median 50\%).
\item 21\%; 78.
\item Eleven.
\item Squarepoint (US) \$2\,567m, Two Sigma Advisers \$735m, AHL Partners \$436m, Renaissance \$307m, AQR \$263m, D. E.
Shaw \$167m, Voleon \$147m, Graham \$140m, Two Sigma Investments \$70.0m, Winton \$69.5m, PDT \$51.9m, CFM \$43.1m,
Aspect \$24.4m.
\item D. E. Shaw, Voleon, Graham, Two Sigma Investments and Winton.
\item Voleon (16\%) and AHL Partners (81\%).
\item Two Sigma Investments and Two Sigma Advisers: one may advise portfolios the other also reports, so their RAUM can
overlap, and their staff are split across entities.
\item Winton, 0.46.
\item Middle half \$58--221 million per employee (median \$119 million); named advisers from \$24 million to \$2.6
billion; top-ten share 21\% and index 78.
\item It divides assets advised through a US entity by that entity's small staff.
\item That most of its staff build and run systems rather than hold advisory roles: a research job there sits in a large
engineering organisation.
\item Revenue, returns, pay, where the rest of the group's staff are, and how the entities relate.
\item As dated ranges of size and research share for the registered entities, read against their filing dates and group
structure; never as productivity or pay.
\end{enumerate}
\end{solution}

\begin{solution}{iq:in:quantitative-hedge-funds-and-systematic-managers:1}
At a systematic fund the researcher's output is a rule that trades without them, judged by out-of-sample performance
and its interaction with the portfolio; at a discretionary fund the researcher supports a portfolio manager's decisions,
and is judged by the manager.
\iqlookfor{who decides the trade, and so what the researcher ships.}
\end{solution}

\begin{solution}{iq:in:quantitative-hedge-funds-and-systematic-managers:2}
Leverage: borrowing, short positions and derivatives add gross exposure without adding capital; a market-neutral book
may hold longs and shorts each larger than its net assets.
\iqlookfor{gross against net.}
\end{solution}

\begin{solution}{iq:in:quantitative-hedge-funds-and-systematic-managers:3}
\$100 million of RAUM per employee (near the median of the file), 40\% in investment work. Before comparing: whether both
rows are the whole firm or one entity, how much RAUM is leverage, the RAUM against private-fund gross assets, the filing
dates, and the business mix.
\iqlookfor{ratios, their scope and their denominators.}
\end{solution}

\begin{solution}{iq:in:quantitative-hedge-funds-and-systematic-managers:4}
To raise money (investors want to understand the process), to hire (researchers want to publish and to see the
culture), and because method is not the edge: published factor research and statistics help sell and recruit without
revealing specific signals.
\iqlookfor{secrecy of signals, openness of method.}
\end{solution}

\begin{solution}{iq:in:quantitative-hedge-funds-and-systematic-managers:5}
No: a top-ten share of 21\% and an index of 78 are low. The answer would change within a strategy (a few firms may
dominate one market or style), by the measure used (net assets, revenue), or if private-fund assets outside the United
States were counted.
\iqlookfor{the market's definition drives the answer.}
\end{solution}

\begin{solution}{iq:in:quantitative-hedge-funds-and-systematic-managers:6}
Usually trend following: futures markets are deep and positions are held for weeks, so costs grow slowly with size;
its limit is the market-impact cost of the largest contracts and crowding in the same trends. The statistical-arbitrage
book turns over fast in many small positions: impact grows with size and its capacity is lower (Book~7, chapter~28).
\iqlookfor{capacity as turnover times impact.}
\end{solution}
